\begin{claim}{F.1}{Proved}{All-set fibre functors recover the fundamental groupoid}
Let $X$ be a $\U$-small space, let $G=\Pi_1X$, and retain the
ordinary category $\mathcal L_0(X)=\Fun(G,\Set_\U)$ together with
its evaluation functors
$\omega_x:\mathcal L_0(X)\to\Set_\U$ for a set of points meeting
every component.  The category whose objects are these functors and
whose arrows are all natural transformations between them is a
groupoid equivalent to $G$.  With every point retained, the
comparison is the specified fully faithful functor
\[
 x\longmapsto\omega_x,
 \qquad \Hom_G(x,y)\longrightarrow
             \operatorname{Nat}(\omega_x,\omega_y).            \tag{F.1.1}
\]
Restricting instead to finite sets recovers the profinite Galois
invariant; it must not be identified with the full groupoid in
(F.1.1).
\end{claim}
\begin{proof}
The covariant representable
$h_x=\Hom_G(x,-)$ belongs to $\mathcal L_0(X)$ and represents
$\omega_x$: the ordinary Yoneda isomorphism is
$\Hom(h_x,L)\cong L(x)$.  Applying Yoneda once more gives
\[
 \operatorname{Nat}(\omega_x,\omega_y)
    \cong\Hom_{\mathcal L_0(X)}(h_y,h_x)
    \cong\Hom_G(x,y).
\]
Under these isomorphisms (F.1.1) sends an arrow to transport along
that arrow; composition is preserved.  Every such transformation is
invertible because every arrow of $G$ is invertible.  Points in
different components have no arrows between their evaluations,
exactly as the displayed formula states.  A subset of points meeting
every component gives an equivalent full subgroupoid.  If $X$ is
empty the family is empty and the reconstructed groupoid is empty,
even though $\mathcal L_0(X)$ is terminal.

The identification of $\mathcal L_0(X)$ as the set-valued local
system category and the corresponding reflector is the checked
CATHOT C.1--C.4.  The above argument reconstructs that groupoid
from its fibres; it does not reprove the CATHOT reflection theorem.
In contrast, SGA~1, Expos\'e~V, Theorem~4.1 explicitly takes a fibre
functor with values in \emph{finite} sets and obtains a topological
automorphism group represented by a pro-object
\citep{EF_SGA1}.  For a connected discrete group $G$ this is its
profinite completion, since finite continuous actions are precisely
the finite actions factoring through finite quotients.  Thus the
finite and all-set statements have different quantifiers and
different conclusions.
\end{proof}

\begin{claim}{F.2}{Proved}{Enhanced linear reconstruction with a specified fibre}
Let $R$ be an $E_\infty$-ring spectrum and $(X,x)$ a pointed connected
$\U$-small space.  Retain the $R$-linear stable $\infty$-category
\[
 \mathcal L_R(X)=\Fun(X,\Mod_R)
\]
and its $R$-linear evaluation functor
$\omega_x:\mathcal L_R(X)\to\Mod_R$, including its coherent
functor structure.  This pair determines the $E_1$-algebra
$R[\Omega_xX]$ up to equivalence of $E_1$-algebras.  For an
ordinary commutative ring $k$ and $R=Hk$, it determines
$C_*(\Omega_xX;k)$ as an associative derived $k$-algebra, with
the prescribed variance convention.  Retaining fibres in every
component gives the corresponding many-object linear path category.
\end{claim}
\begin{proof}
The functor $\omega_x$ has a left adjoint $F_x$, unique with its
adjunction data up to contractible choice.  Left Kan extension gives
the explicit formula
\[
 F_x(M)(y)=R[\Map_X(x,y)]\otimes_R M.
                                                               \tag{F.2.1}
\]
Here $\Map_X(x,y)$ is the space of paths from $x$ to $y$, and
$R[T]=R\otimes\Sigma^\infty_+T$.  Thus the monad
$T_x=\omega_xF_x$ is tensoring with $R[\Omega_xX]$.
Its unit comes from the constant path and its multiplication from
concatenation, with the order fixed by the unit and counit of this
adjunction.  Evaluating the $R$-linear monad at the free rank-one
module recovers
\[
 T_x(R),\qquad
 R\longrightarrow T_x(R),\qquad
 T_x(R)\otimes_R T_x(R)\longrightarrow T_x(R),                 \tag{F.2.2}
\]
together with their coherent associativity and unit data.  The
linear functor theorem
\citep[Theorem 4.8.4.1]{EF_LurieHA} identifies $R$-linear
colimit-preserving endofunctors of $\Mod_R$ with tensoring by
$R$-modules, so (F.2.2) is a reconstruction of the algebra itself.
It does not reconstruct merely a Morita class.

The equivalence
$\mathcal L_R(X)\simeq\Mod_{R[\Omega_xX]}$ over $\Mod_R$
is precisely the monadic equivalence proved in E.2: evaluation is
conservative on a connected groupoid and preserves geometric
realizations.  Its proof therefore includes the hypotheses of the
Barr--Beck theorem, rather than inferring the assertion from a
comparison of homotopy categories.
For $R=Hk$, the symmetric monoidal chain construction identifies
$Hk\otimes\Sigma^\infty_+\Omega_xX$ with the chain algebra of
the based loop space.  In the presheaf convention the opposite
action is transported by path inversion; with right modules one
instead obtains the opposite algebra.  Either convention is valid,
but the comparison map is fixed after the convention is chosen.

With more than one point, enriched adjunction gives the mapping
$R$-module
\[
 \underline{\Hom}_{\mathcal L_R(X)}(F_y(R),F_x(R))
                 \simeq R[\Map_X(x,y)].                       \tag{F.2.3}
\]
Composition on the left, with this reversal of the representables,
is path composition on the right.  Formula (F.2.3), component by
component, reconstructs the $R$-linearized path category.  Between
different components the path space is empty and its linearization
is zero.  No basepoint is chosen in an empty component.
\end{proof}

\begin{claim}{F.3}{Proved}{Spaces are recovered from enhanced space-valued local systems}
For every $\U$-small space $X$, the $\infty$-topos
$\mathcal L_{\mathcal S}(X)=\Fun(X,\mathcal S_\U)$ has shape
the constant pro-space $X$.  Retaining its evaluation functors
$\omega_x$ and their full spaces of natural transformations
recovers $X$ as an $\infty$-groupoid.  Retaining only the
homotopy category still suffices if, in addition, the specified
constant functor
$\Ho\mathcal S\to\Ho\mathcal L_{\mathcal S}(X)$ is retained:
in that case one recovers the object $X$ of $\Hot$ by Yoneda.
\end{claim}
\begin{proof}
Straightening identifies $\Fun(X,\mathcal S)$ with the slice
$\infty$-topos $\mathcal S_{/X}$.  The inverse-image part
$\Delta_X:\mathcal S\to\mathcal L_{\mathcal S}(X)$ of its
geometric morphism to spaces is the constant functor.  Its global
sections satisfy the canonical equivalence
\[
 \Map_{\mathcal L_{\mathcal S}(X)}(1,\Delta_XK)
                         \simeq\Map_{\mathcal S}(X,K).         \tag{F.3.1}
\]
By definition the left side, as a functor of $K$, is corepresented
by the shape.  Equation (F.3.1) identifies it with the constant
pro-object represented by $X$.  This is also the classifying-topos
case of the shape construction in
\citet[Section 2]{EF_Hoyois2018}; it is a known reconstruction,
not a priority claim.

For the fibre-functor formulation, the covariant representable
$r_x(y)=\Map_X(x,y)$ corepresents $\omega_x$.  The enriched
Yoneda equivalence yields
\[
 \Map_{\Fun(\mathcal L_{\mathcal S}(X),\mathcal S)}
       (\omega_x,\omega_y)
 \simeq\Map_{\mathcal L_{\mathcal S}(X)}(r_y,r_x)
 \simeq\Map_X(x,y).                                          \tag{F.3.2}
\]
These equivalences respect composition coherently.  They identify
$X$ with the full $\infty$-subcategory on its specified evaluation
functors; every morphism there is an equivalence.  If only a family
meeting every component is retained, the resulting full subgroupoid
is equivalent to $X$.

Finally taking components of (F.3.1) gives
$\Hom(1,\Delta_XK)=[X,K]$, naturally in $K\in\Hot$.
The ordinary Yoneda lemma recovers the representing object $X$ in
$\Hot$ from these hom-sets.  The extra datum of the constant
functor is essential to this particular proof: it supplies the
specified functor whose representability is being used.  None of
these arguments identifies an arbitrary equivalence of bare
homotopy categories with a geometric equivalence.  For $X=\varnothing$
the topos is degenerate, (F.3.1) is the contractible space
$\Map(\varnothing,K)$, and the same representability assertion
recovers the empty space correctly.
\end{proof}

\begin{claim}{F.4}{Known}{The form applicable to a locally connected topos}
Let $\mathcal E$ be a locally connected Grothendieck topos, and
let $\Pi_1\mathcal E$ be its fundamental pro-groupoid, equivalently
the 1-truncation of the shape of its associated $\infty$-topos.
There is a canonical equivalence
\[
 \Fun(\Pi_1\mathcal E,\Set)
             \simeq\LC_{\Set}(\mathcal E).                   \tag{F.4.1}
\]
For a representing cofiltered diagram $(G_i)$, the left side means
the filtered colimit of the categories $\Fun(G_i,\Set)$ in the
sense of representations of a pro-groupoid; it is not the category
of actions of the abstract inverse-limit group with its topology
forgotten.  The full subcategory of filtered colimits of locally
constant objects is the corresponding classifying topos.  In
particular, if $\mathcal E$ is generated under colimits by its
locally constant objects, the canonical map
$\mathcal E\to\beta(\Pi_1\mathcal E)$ is an equivalence.
\end{claim}
\begin{proof}[Source, hypotheses, and specialization]
Equation (F.4.1) is
\citet[Theorem 3.13, case $n=0$]{EF_Hoyois2018}, with
0-truncated coefficients.  The proof uses local connectedness to
make the constant-object functor preserve products on sets, then
descends the resulting mapping-space identities along torsors.
The description of the filtered-colimit completion is part of that
same theorem.  The final assertion is
\citet[Corollary 3.20, case $n=1$]{EF_Hoyois2018}, which permits
several components and does not assume a point.  For an
$\infty$-topos locally $n$-connected, the same theorem identifies
locally constant $n$-truncated objects with
$\Fun(\Pi_\infty\mathcal E,\mathcal S_{\le n})$.
Without any connectedness hypothesis, its finite version is
\citet[Theorem 4.3]{EF_Hoyois2018}, with finite homotopy types as
coefficients.

This also specifies the reconstruction data needed by the
multi-Galois application.  Retain the category of locally constant
objects with finite limits, and the families which are jointly
effective epimorphic in $\mathcal E$.  For each small groupoid $G$
these data determine its principal $G$-torsors: they consist of
objects indexed by the objects of $G$, transport isomorphisms, the
principal-action identities, and the joint covering condition.
The resulting groupoid-valued functor of $G$ is the torsor functor
corepresented by $\Pi_1\mathcal E$.  Yoneda in pro-groupoids
determines this pro-object up to equivalence.  If a family of points
is jointly conservative on $\mathcal E$, their
fibre functors identify the covering families as those which are
jointly surjective on every fibre.  Without this conservativity
assumption one must retain the covering condition itself; a
statement about unspecified point functors would have omitted data.

The attached \emph{Longue Marche}, Sections 1--4, uses the stronger
locally simply connected setting when asserting that locally
constant objects already form a multigaloisian topos
\citep[Section 2, Theorem 2.1 and Definition 2.2]{EF_LongueMarche}.
For merely locally connected $\mathcal E$, the distinction in
(F.4.1) between split representations and their filtered-colimit
completion is retained.  We do not replace it by an unqualified
assertion that $\LC(\mathcal E)$ is always a topos.
\end{proof}

\begin{claim}{F.5}{Known}{Kaledin's groupoid enhancement already reconstructs spaces}
For a small space $X$ consider the family
\[
              J\longmapsto\Pi_1\Map(BJ,X).                   \tag{F.5.1}
\]
Kaledin's enhanced-groupoid representability theorem identifies
the appropriate small families satisfying his axioms with homotopy
types.  More precisely, in the attached arXiv:2409.17489v3,
Definition 6.1.1.3 and Corollary 6.1.2.5 give the groupoid case;
Theorem 7.3.1.3 is the more general representability theorem for
enhanced categories.  Thus the reconstruction of a space from
(F.5.1), with all the structure required there, is already recorded
in the literature.
\end{claim}
\begin{proof}[Exact comparison with the stated theorem]
Kaledin's $\mathrm{Pos}$ is the category of partially ordered
sets of finite chain dimension, not the category of finite posets.
Its objects may have arbitrary permitted cardinality.  His
$\mathrm{Pos}^{+}$ consists of left-bounded partially ordered
sets.  A small enhanced groupoid is a family of groupoids over
$\mathrm{Pos}$ satisfying additivity, semiexactness, and excision
as specified in Definition 6.1.1.1.  A semicanonical extension to
$\mathrm{Pos}^{+}$ also has the stated interval-invariance and
semicontinuity properties.  The small or $\kappa$-bounded
conditions, for an uncountable regular $\kappa$, are part of the
claim and cannot be omitted.

The comparison of these families with families on simplicial sets
is Proposition 6.1.2.3.  Corollary 6.1.2.5 identifies their category
with the localization of simplicial sets at weak equivalences,
and supplies uniqueness of extension with the precise equivalence
and isomorphism ambiguities stated there.  In the representable
case its nerve pullback is exactly (F.5.1), with derived mapping
spaces.  Theorem 7.3.1.3 proves that every restricted enhanced
category is equivalent over $\mathrm{Pos}$ to one obtained from
a complete fibrant Segal space; it also represents enhanced
functors by maps and identifies homotopic maps with isomorphic
enhanced functors \citep{EF_Kaledin2026}.

These are statements of the opened preprint.  Its arXiv record
specifies that Chapters 6 and 7 remain preliminary, whereas
Chapters 1--5 are intended for publication.  We therefore do not
describe the enhanced-groupoid theorem as a checked consequence of
a refereed 2026 journal article.  The space reconstruction F.3 was
proved independently above and does not depend on the status of
that preprint.
\end{proof}

\begin{claim}{F.6}{Proved}{The relation between the coefficient tower and varying shapes}
The finite-type tower in E.5 and the family (F.5.1) both determine
$X$ up to equivalence, but their parameters and their truncations
are different.  The tower raises the coefficient category at fixed
$X$.  Kaledin fixes the target of the reported values to ordinary
groupoids and varies the source shape $J$, retaining all coherence
between those values.  In general
\[
 \Pi_1\Map(BJ,X)\not\simeq
             \Map(BJ,\Pi_1X),                                \tag{F.6.1}
\]
even on components.  The explicit comparison from enhanced
space-valued local systems is obtained by first forming the full
$\infty$-groupoid of designated fibre functors in F.3 and then
taking $\Pi_1\Fun(BJ,-)$ of that groupoid.
\end{claim}
\begin{proof}
For the coefficient tower, E.5 says that a specified map inducing
equivalences at all stages induces isomorphisms on all homotopy
groups.  For varying shapes, F.5 gives the known representability
statement.  Formula (F.3.2) identifies the full fibre-functor
$\infty$-groupoid with $X$, compatibly with composition, so the
last construction gives (F.5.1) naturally in $J$.

For (F.6.1), take $X=K(\Z,2)$ and take $J$ to be the face poset
of a finite triangulation of $S^2$.  Then $BJ\simeq S^2$ and
\[
 \pi_0\Pi_1\Map(BJ,X)
       =[S^2,K(\Z,2)]\cong\Z,
 \qquad \Map(BJ,\Pi_1X)\simeq *.
\]
The left side contains degree-two information even though each
reported value is an ordinary groupoid.  Thus one cannot identify
Kaledin's diagram with a diagram built only from the unenhanced
groupoid $\Pi_1X$.  This example involves an infinite homotopy
type as target; the finite poset itself is the nonempty proper
subsets of a four-element set ordered by inclusion, whose nerve
is the barycentric subdivision of the boundary of a tetrahedron.
Its exact face counts and homology are verified in
\texttt{code/ef\_checks.py} and its output.
\end{proof}

\begin{claim}{F.7}{Proved}{A precise recognition question for a coefficient theory}
For a complete coefficient $\infty$-category $\mathcal C$,
the condition $W_{\mathcal C}^{\mathrm{lc}}=W_\infty$ means
that the contravariant functor
$X\mapsto\Fun(X,\mathcal C)$ reflects equivalences of
\emph{specified maps} of spaces.  Equivalently, the pair of
test spaces $K_0(\mathcal C),K_1(\mathcal C)$ of E.1 detects
all weak equivalences.  In the cases computed in Part C, this
holds for spaces, integral derived coefficients, and spectra;
it fails for sets, groupoids, every fixed finite truncation level,
rational derived coefficients, and mod-$p$ derived coefficients.
\end{claim}
\begin{proof}
The equivalence of the two formulations is E.1.  The positive
and negative computations are those of Part C; F.1--F.3
describe their reconstructed invariants with fibres retained.
For finite truncations the projection of a nontrivial
$K(A,r+2)$ to a point passes the $\mathcal S_{\le r}$ local
system test but is not a weak equivalence.  For rational
coefficients use $K(\Z/p,2)\to *$.  For mod-$p$ coefficients
use $K(\Z/q,2)\to *$ with a different prime $q$; transfer
and the path fibration give the needed mod-$p$ acyclicity.
All spaces in these examples are connected, so no ambiguity
about component detection is involved.  For a general ring
or an $E$-local stable coefficient theory, the exact loop-algebra
criteria and the qualifications in Part C remain in force.
\end{proof}

\begin{claim}{F.8}{Open}{The reconstruction problem with only ordinary fibre functors}
If the input in Part F is restricted literally to the ordinary
category $\LC_{\D_R}(X)$ and ordinary functors
$\LC_{\D_R}(X)\to\D_R(e)$, with no enhanced mapping objects,
no coherent adjunction data, and no specified family of shifted
categories, F.2 does not establish reconstruction of the
$E_1$-algebra $C_*(\Omega X;k)$.  We leave that unenhanced
assertion unresolved.  Likewise F.3 without either its
$\infty$-topos structure or its specified constant functor
does not establish the corresponding bare-category assertion.
\end{claim}
\begin{proof}[The precise missing comparison]
The proof of F.2 constructs an algebra from a coherent monad.
The ordinary adjunction on homotopy categories supplies a
monad in ordinary categories, but no comparison reconstructing
all of the higher multiplication coherences has been supplied.
An invalid passage from the latter monad to an $E_1$-algebra
would be a gap, not a counterexample.  A smallest useful
test is the class of pointed simply connected finite CW
complexes with rational coefficients: does an equivalence
of their ordinary local-system categories, compatible with
the specified ordinary evaluation functors, imply equivalence
of the rational loop-chain $E_1$-algebras?  Neither a proof
nor a counterexample to this exact implication is given here.
The enhanced assertion and the constant-functor space assertion
have been proved with their extra data explicitly named.
\end{proof}
